%%
%% fall-lecture22.tex
%% 
%% Made by Alex Nelson <pqnelson@gmail.com>
%% Login   <alex@lisp>
%% 
%% Started on  2026-11-17T11:51:23-0800
%% Last update 2026-11-17T11:51:23-0800
%% 

\lecture[Integrable functions]{}

\begin{node}
Let $(X,\mathcal{A},\mu)$ be a measure space. We will define
integrable functions on $X$ in three steps.
\end{node}

\begin{definition}
Let $(X,\mathcal{A},\mu)$ be a measure space.
\begin{enumerate}
\item 
A \define{Simple Function} is a function $f\colon X\to\RR$ such that
$f=\sum_{j} c_{j}\mathbf{1}_{A_{j}}$ where $A_{j}\in\mathcal{A}$ are
disjoint [or, the pairwise intersection $A_{i}\cap A_{j}$ is a nullset
  when $i\neq j$] [and the $A_{j}$ cover $X$], and
$c_{j}\in[0,\infty]$ is a non-negative extended real number.
\item When $f=\sum_{j}c_{j}\mathbf{1}_{A_{j}}$ is a simple function,
  we define the \define{Integral} of $f$ to be the extended real
  number $\int f\,\D\mu=\sum_{j}c_{j}\mu(A_{j})$.
\end{enumerate}
\end{definition}

\begin{definition}
Let $(X,\mathcal{A},\mu)$ be a measure space.
Let $f\colon X\to\RR$ be a measurable function~(\S\ref{defn:lec21:measurable-function}).
We define the \define{Integral} of $f$ to be the extended real number
$\int f\,\D\mu:=\sup_{s}\{\int s\,\D\mu\mid 0\leq s\leq f,s\mbox{ is simple}\}$.
\end{definition}

\begin{definition}
Let $f\colon X\to\RR$ be a measurable function.
\begin{enumerate}
\item We say $f$ is
\define{Integrable} if $\int|f|\,\D\mu<\infty$ is finite.
\item We define the \define{Integral} of $f$ to be $\int f\,\D\mu=\int f^{+}\,\D\mu-\int f^{-}\,\D\mu$
  where $f^{+}(x)=\max(0,f(x))$ and $f^{-}(x)=\max(0,-f(x))$ and $f=f^{+}-f^{-}$.
  [This just extends the notion of the integrable of a non-negative
    measurable function by linearity.]
\end{enumerate}
\end{definition}

\begin{node}
In probability, we write $\int X\,\D\Pr=\mathbb{E}[X]$ --- it's the
expected value of the random variable.
\end{node}

\begin{notation}
When writing the integral $\int f\,\D\mu$, we will suppress the
$\D\mu$ measure when it is clear from context.
\end{notation}

\begin{proposition}
Let $f\geq0$ be measurable. Then $\int f=0$ if and only if $f=0$
almost everywhere.
\end{proposition}

\begin{proof}
\backwardproof\ Assume $f=0$ almost everywhere. Let $0\leq s\leq f$ be
a simple function approximating $f$. Then $s=0$ almost everywhere, so
$\int s=0$. Then $\int f=\sup\{\int s\mid 0\leq s\leq f, s\mbox{ simple}\}=\sup\{0\}=0$.

\forwardproof\ Assume $\int f=0$. Write $E_{n}=\{f>1/n\}$. Suppose for
contradiction that $f$ is nonzero on a set of positive measure. Then
there exists an $n$ such that $\mu(E_{n})>0$. So $f>\frac{1}{n}\mathbf{1}_{E_{n}}$,
and $\frac{1}{n}\mathbf{1}_{E_{n}}$ is a positive simple
function. Therefore
\begin{equation}
\int>\frac{1}{n}\mu(E_{n})\neq0,
\end{equation}
which is a contradiction. Therefore $f$ is zero almost everywhere.
\end{proof}

\begin{definition}
We say two measurable functions $f$ and $g$ are \define{Equivalent}
written $f\sim g$ if $f=g$ almost everywhere (i.e., $\int|f-g|=0$).
\end{definition}

\begin{definition}
We may define the \define{Lebesgue Space} $L^{1}(\mu)=L^{1}(X,\mathcal{A},\mu)$
is the equivalence classes of integrable functions.

We also have the \define{$L^{1}$ Norm} $\|f\|_{L^{1}}:=\int|f|$ which
is a function $\|-\|_{L^{1}}\colon L^{1}(\mu)\to[0,\infty)$. (We have
  to prove it satisfies the properties of the norm.)
\end{definition}

\begin{theorem}
\begin{enumerate}
\item $L^{1}(\mu)$ is a normed vector space (with respect to the $L^{1}$-norm)
\item The map $f\mapsto\int f\,\D\mu$ is a continuous linear
  functional of norm 1.
\end{enumerate}
\end{theorem}

\begin{remark}
\begin{enumerate}
\item Also, $L^{1}$ is complete---so $L^{1}$ is a Banach space.
\item And also, simple functions are dense in $L^{1}$.
\end{enumerate}
\end{remark}

\begin{claim}[Approximation by simple functions]
Let $f\geq0$ be measurable and monotone. Then there exists an
increasing monotone sequence of simple functions $0\leq s_{1}\leq s_{2}\leq\dots\uparrow f$
pointwise.
\end{claim}

\begin{proof}
For any $n\in\NN$, we will construct an $\varepsilon$-approximation of
$f$ where $\varepsilon=2^{-n}$. Define
\begin{equation}
E_{k}=\{\frac{k}{2^{n}}\leq f<\frac{k+1}{2^{n}}\}
\end{equation}
dyadic intervals of width $\varepsilon$, set
$s_{n}:=\sum\frac{k}{2^{n}}\mathbf{1}_{E_{k}}$.
This gives us a monotone increasing sequence $(s_{n})$ which converges
to $f$.
\end{proof}

\subsection{Convergence theorems}

\begin{monotone-convergence-theorem}[Levy]
If we have a sequence of functions $0\leq f_{n}$ which pointwise
increases towards $f$, then $\int f_{n}\nearrow\int f$ converges monotone.
\end{monotone-convergence-theorem}

\begin{proof}
Let us set
\begin{equation}
a := \lim_{n\to\infty}\int f_{n}.
\end{equation}
Then $a\leq\int f$ since each $f_{n}\leq f$ (so $\int f_{n}\leq\int f$
for each $n$). We want to show $\lim\int f_{n}\geq\int f$ [so then we
obtain equality].
We will fix some $0\leq s\leq f$ simple function, and fix some $0<c<1$.
[``Fix'' meaning that we will eventually take the limit $s\to f$
  (resp., $c\to 1$).]
Let
\begin{equation}
E_{n}:=\{x\mid f_{n}(x)\geq cs(x)\}.
\end{equation}
The main observation is that the sets $E_{n}$ are monotone increasing towards
$X$. So $\lim_{n\to\infty}E_{n}=\bigcup_{n}E_{n}$. Now, we have
\begin{equation}
\int_{X}f\geq\int_{E_{n}}f_{n}\geq c\int_{E_{n}}s,
\end{equation}
and
\begin{equation}
c\int_{E_{n}}s\xrightarrow{n\to\infty}c\int_{X}s\quad\mbox{converges}.
\end{equation}
We use the fact that measures are continuous on monotone sequences (if
$E_{n}\uparrow E_{\infty}$, then $\mu(E_{n})\xrightarrow{n\to\infty}\mu(E_{\infty})$
[``continuity from below'']).
We apply this continuity theorem to the function $E\mapsto\int_{E}s$
[this is a very important operation, we'll discuss with the
  Radon--Nikodyn theorem]. Then
\begin{equation}
a=\lim_{n\to\infty}\int_{X}f_{n}\geq c\int_{X}s,
\end{equation}
then take $c\to1$ limit, and then take the supremum over $s$. Then we
obtain
$a\geq c\int_{X}s$. Hence equality.
\end{proof}

\begin{corollary}
Suppose all $f_{n}\geq0$. Then $\sum(\int f_{n})=\int(\sum f_{n})$.
\end{corollary}

\begin{corollary}
If $\sum\int|f_{n}|<\infty$ is finite, then $\sum f_{n}$ absolutely
converges almost everywhere.
\end{corollary}